% 前言
\renewcommand{\codeInset}{18bp}
C++ 是一种通用编程语言，最初的设计偏向嵌入式编程与系统编程。多年来，C++ 不断演进，如今被用于开发众多不同领域的软件。凭借其多面性与健壮性，C++ 是开启编程之旅的绝佳语言。本书涵盖用 C++ 构建的一系列精彩项目，展示如何在不同场景中运用这门语言。在开发这些项目的过程中，你不仅会学到语言的各个构造，还会学到如何用 C++ 满足自己的软件需求。

本书研究一组用 C++ 编写的应用程序，涵盖抽象数据类型、图书馆管理系统、图形应用、游戏以及领域特定语言（DSL）。

\section*{本书面向的读者}

本书面向想用 C++ 开发软件的开发者。具备基本的编程经验会是额外的优势。

\section*{本书涵盖的内容}

第 1 章「初识 C++」，介绍 C++ 的面向对象编程。我们先研究一个掷骰子的简单程序：编写代码，然后编译、链接并执行。接着构造一个简单的面向对象层次结构，涉及指针与动态绑定。最后创建两个简单的抽象数据类型：栈与队列。栈是自底向上排列的一组值，只有栈顶的值可访问；队列则是传统意义上的队列，在前端读取值、在后端加入值。

第 2 章「数据结构与算法」，在上章所学的基础上展开，特别是列表与集合抽象数据类型。我们还引入模板与运算符重载，并研究线性搜索与二分搜索算法，以及插入排序、选择排序、冒泡排序、归并排序与快速排序算法。

第 3 章「构建图书馆管理系统」，帮助你开发一个真实系统：一个由图书和读者组成的图书馆管理系统。图书记录哪些读者借阅和预约了它，读者则记录自己借阅和预约了哪些图书。

第 4 章「基于指针的图书馆管理系统」，进一步发展图书馆管理系统。上一章中每本图书和每个读者都用整数编号标识；而在本章中，我们改用指针。每本图书持有指向借阅或预约它的读者的指针，每个读者则持有指向自己借阅或预约的图书的指针。

第 5 章「Qt 图形应用」，深入三个用 Qt 图形库开发的图形应用：一个带时针、分针和秒针的模拟时钟，一个用不同颜色绘制直线、矩形和椭圆的画图程序，以及一个让用户输入并编辑文本的编辑器。我们将学习在 Qt 库中处理窗口与部件、菜单与工具栏，还会学习如何绘制图形、书写文本，以及如何捕获鼠标和键盘输入。

第 6 章「增强 Qt 图形应用」，进一步发展三个图形应用：模拟时钟、画图程序和编辑器。我们为钟面添加数字，为画图程序添加移动、修改与剪切粘贴图形的功能，并为编辑器添加更改字体与文本对齐方式的功能。

第 7 章「游戏」，介绍基础的游戏开发。我们在本章用 Qt 库开发奥赛罗棋和圈叉棋两种游戏。在奥赛罗棋中，双方轮流在棋盘上放置黑色和白色的棋子，以包围对手的棋子；在圈叉棋中，双方轮流放置圆圈和叉号，以将五个棋子连成一行。

第 8 章「计算机对弈」，让计算机与人类玩家对弈。在奥赛罗棋中，计算机试图放置棋子以包围尽可能多的对手棋子；在圈叉棋中，计算机试图放置棋子使自己的五个棋子连成一行，同时阻止对手连成一行。

第 9 章「领域特定语言」，教你开发领域特定语言（DSL），即面向特定领域的语言。更具体地说，我们开发一种在 Qt 部件上书写图形对象的语言。我们先用文法形式化地定义这门语言，然后编写识别有意义字符序列的扫描器、检查源代码是否符合文法的解析器，以及显示图形对象的查看器。

第 10 章「高级领域特定语言」，从多个方面改进我们的领域特定语言：添加改变程序流程的选择与迭代，添加可在程序执行期间赋值的变量，并添加带参数与返回值的函数。

\section*{充分利用本书}

本书面向每一位读者，从初学者到较为熟练的 C++ 程序员皆可阅读。不过，具备一些 C++ 经验会更有帮助。

本书的示例在 Visual Studio 与 Qt Creator 中开发。

\subsection*{下载示例代码文件}

你可以从 www.packtpub.com 的个人账户下载本书的示例代码文件。如果你是在别处购买的本书，可以访问 www.packtpub.com/support 并注册，以便将文件直接发送到你的邮箱。

你可以按以下步骤下载代码文件：

\begin{enumerate}
\item 登录或注册 www.packtpub.com。
\item 选择 SUPPORT 标签页。
\item 点击 Code Downloads \& Errata。
\item 在搜索框输入书名，并按照屏幕上的指示操作。
\end{enumerate}

文件下载完成后，请确保用以下软件的最新版本解压文件夹：

\begin{itemize}
\item Windows：WinRAR/7-Zip
\item Mac：Zipeg/iZip/UnRarX
\item Linux：7-Zip/PeaZip
\end{itemize}

本书的代码包也托管在 GitHub：https://github.com/PacktPublishing/CPP17-By-Example。我们丰富的图书和视频目录中的其他代码包也发布在 https://github.com/PacktPublishing/。欢迎查看！

\subsection*{下载彩色图片}

我们还提供一个 PDF 文件，收录本书中使用的截图/示意图的彩色版本，可在此下载：\allowbreak{\fontsize{9.5bp}{12bp}\selectfont https://www.packtpub.com/sites/default/files/downloads/CPP17ByExample\_ColorImages.pdf}。

\subsection*{本书中使用的约定}

本书使用了若干文本约定。

\texttt{CodeInText}：表示正文中的代码字词。例如：「我们继续看一个类层次结构，其中 \texttt{Person} 是基类，\texttt{Student} 与 \texttt{Employee} 是它的派生类：」

代码块按如下方式排版：

\begin{codebox}
class Person {
   public:
     Person(string name);
     virtual void print();
   private:
     string m_name;
};
\end{codebox}

当我们希望你注意代码块中的特定部分时，相关的行或条目会以加粗显示（原书中 \texttt{print} 一行为加粗等字体；本译稿统一以正体呈现）：

\begin{codebox}
class Person {
   public:
     Person(string name);
     virtual void print();
   private:
     string m_name;
};
\end{codebox}

\textbf{Bold}：表示新术语、重要词语，或者你在屏幕上看到的字词。例如，菜单或对话框中的词语在正文中显示如下：「在第一个对话框中，我们只需按下 Next 按钮。」

\begin{myNotic}{p017_i00}{重要提示或注意事项以此样式出现。}
\end{myNotic}

\begin{myNotic}{p017_i01}{提示与技巧以此样式出现。}
\end{myNotic}

\section{联系我们}

我们始终欢迎读者的反馈。

\textbf{一般反馈}：请发邮件至 feedback@packtpub.com，并在邮件主题中注明书名。如果你对本书的任何方面有疑问，请发邮件至 questions@packtpub.com。

\textbf{勘误}：尽管我们在内容的准确性上已经尽最大努力，错误仍在所难免。如果你发现书中有错误，恳请告知。请访问 www.packtpub.com/submit-errata，选择你的书籍，点击勘误提交表单链接，并输入详情。

\textbf{盗版}：如果你在互联网上遇到任何形式的非法复制本，恳请告知我们地址或网站名称，并附上相关链接，联系 copyright@packtpub.com。

\textbf{成为作者}：如果你精通某个主题并有兴趣写书或参与写作，请访问 authors.packtpub.com。

\subsection{留下一条书评}

请留下一条书评。读完并使用本书之后，不妨在你购买本书的站点上留下一条书评。这样，潜在读者可以看到并借鉴你的客观意见，做出购买决定；我们 Packt 也能了解你对产品的看法，我们的作者也能看到你对其著作的反馈。谢谢！

要了解更多关于 Packt 的信息，请访问 packtpub.com。
